\chapter{DÜŞÜK GÜÇ ANALİZİ VE YENİ ANAHTARLAMA YÖNTEMİ}
\label{chap:anahtarlama}

Bölüm~\ref{chap:sistem}, T-tipi faz kolunun gerilim durumlarını, akım yollarını, düşük güçteki senkron-buck eşdeğerini, LCL filtre modelini ve doğrulanmış benzetim modelini ortak taban olarak kurmuştur. Bölüm~\ref{chap:anahtarlama}, aynı taban üzerinde tezin ilk katkısını, nötr anahtar susturma yöntemini ele almaktadır. Bölümde ters akımın oluşumu incelenmekte, klasik anahtarlama (S1) ile önerilen yöntem (S2) karşılaştırılmaktadır. Karşılaştırma iki kanıt türüne dayanmaktadır: deney özet tablosu ve deneyin aktif ve reaktif gücüne kalibre edilmiş benzetim modeli. İki kanıt türü ayrı raporlanmıştır.

T-tipi dönüştürücünün düşük güç bölgesindeki davranışını aktarılan aktif gücün azalması tek başına açıklamamaktadır. Düşük güçte anahtarlama dalgalanması ortalama endüktör akımına göre büyük kalmakta ve akım bir anahtarlama periyodu içinde sıfıra ulaşmaktadır. Akım sıfıra ulaştığında klasik senkron anahtarlama ters yönde akıma izin verebilir. Yararlı çıkış gücüne katkı sağlamayan ters akım, yarı iletkenlerin ve filtre elemanlarının RMS akımını artırarak hafif yük verimini düşürmektedir. Ters akım yolunun kesilmesi ise süreksiz iletime ve akım dalga biçiminde özellikle üçüncü harmonik bileşenin büyümesine neden olabilir. Düşük güçteki temel tasarım problemi, verim kazancı ile harmonik bozulma arasındaki ödünleşimin yönetilmesidir.

Düşük güç bölgesi, laboratuvar test noktasının ötesinde gerçek bir işletme koşuludur. Düşük güç koşulu, fotovoltaik eviricilerde gün doğumu ve gün batımı çevresinde, enerji depolama sistemlerinde ise düşük akımlı şarj veya deşarj komutlarında ortaya çıkmaktadır. Fotovoltaik üretim profilleri ile saha ölçümlerinde kısmi yük süresinin ve düşük akımdaki THD artışının gösterilmesi, düşük güç incelemesinin işletme gerekçesini desteklemektedir \citep{Mondol2007,AbRahman2020,Chicco2009}. Enerji depolama uygulamalarında çift yönlü güç komutunun geniş bir aralıkta değişebilmesi de dönüştürücünün yalnız anma noktasında değerlendirilmemesini gerektirmektedir \citep{Aneke2016}.

\section{Ters Akım Dolaşım Mekanizması}
\label{sec:ters-akim-mekanizmasi}

Şebekeye verilen faz akımının temel bileşeni $i_g(t)=\hat I_g\sin(\omega_0t)$ olarak ele alınsın. Düşük güçte $\hat I_g$ küçülürken anahtarlama dalgalanmasının genliği aynı oranda küçülmemektedir. Senkron-buck eşdeğerine göre pozitif yarım periyotta faz kolu $P$ ve $O$, negatif yarım periyotta $N$ ve $O$ durumları arasında anahtarlanmaktadır. Ortalama endüktör akımı dalgalanmanın yarısından küçük kaldığında (Eşitlik~\eqref{eq:ters-akim-esigi}) akım anahtarlama periyodu içinde sıfırı geçmektedir. Sistem parametre değerleriyle akım dalgalanması en çok yaklaşık \SI{9}{\ampere} tepe--tepe, sınır akımı (Eşitlik~\eqref{eq:sinir-akimi}) en çok yaklaşık \SI{4.5}{\ampere}'dir (Alt Bölüm~\ref{sec:dusuk-akim-modelleme}). Deney çalışma noktasının yaklaşık \SI{3.3}{\ampere} tepe akımı sınır akımının altında kaldığından klasik anahtarlama düşük güçte temel periyodun büyük bölümünde zorlanmış sürekli iletim bölgesinde çalışmaktadır.

T-tipi faz kolunun nötr durumu, iki yarım periyotta senkron alçaltıcı dönüştürücünün serbest dolaşım durumuna karşılık gelmektedir. Faz akımı referans yönünde azalarak sıfıra ulaştığında nötr kol iletimde kalırsa LCL filtre kapasitöründe ve şebeke eşdeğerinde bulunan enerji ters yönde akım oluşturmaktadır. Akımın işaret değiştirdiği aralıkta enerji şebekeden veya filtre kapasitöründen nötr noktaya doğru taşınmaktadır. Aktarılan ortalama aktif gücü artırmayan dolaşım akımı, iletim ve bakır kayıpları açısından ek bir yüktür.

Senkron alçaltıcı dönüştürücüde tamamlayıcı sürme, serbest dolaşım akımını diyot yerine MOSFET kanalından geçirerek iletim kaybını azaltmaktadır. İncelenen T-tipi yapının nötr kolu IGBT'lerden oluştuğundan nötr yoldaki akım her zaman bir IGBT ile bir diyottan geçmekte ve kanal iletimi kazancı oluşmamaktadır. Dolayısıyla nötr kolda tamamlayıcı sürmenin düşük güçteki başlıca etkisi, ters dolaşım yolunu açık tutmasıdır.

Bir temel periyot $T_0$, faz kolu endüktör akımı $i_L(t)$, istenen enerji aktarım yönünü veren akım referansı $i^{\ast}(t)$ ve işaret fonksiyonu $\operatorname{sgn}(\cdot)$ ile gösterilsin. Referans yönüne ters akımın akım--zaman büyüklüğünü veren ters akım yükü $Q_{\mathrm{rev}}$
\begin{equation}
  Q_{\mathrm{rev}}=\int_0^{T_0}\max\!\left[-i_L(t)\,\operatorname{sgn}\!\left(i^{\ast}(t)\right),0\right]\,\mathrm{d}t
  \label{eq:ters-akim-yuku}
\end{equation}
ile tanımlanabilir. Eşitlik~\eqref{eq:ters-akim-yuku}, akımın mutlak yönü yerine referansa göre ters bileşeni seçmektedir. Ters akım yükünün birimi amper--saniyedir. Ters akım yükü, kayıp enerjisi ya da güç yerine farklı anahtarlama yöntemlerindeki dolaşımı karşılaştırmak için kullanılmaktadır. Parçalı doğrusal yarı iletken modelinde sabit iletim gerilimi $V_{\mathrm{on}}$, eğim direnci $R_{\mathrm{on}}$, iletimdeki mutlak akım ortalaması $I_{\mathrm{avg,abs}}$ ve RMS akım $I_{\mathrm{rms}}$ ile basitleştirilmiş iletim kaybı
\begin{equation}
  P_{\mathrm{cond}}\simeq V_{\mathrm{on}} I_{\mathrm{avg,abs}}+R_{\mathrm{on}}I_{\mathrm{rms}}^2
  \label{eq:iletim-kaybi-yaklasik}
\end{equation}
ile yaklaştırılabilir \citep{Chen2010,Schweizer2013}. Eşitlik~\eqref{eq:iletim-kaybi-yaklasik} anahtarlama enerjisini ve sıcaklık bağımlılığını içermeyen ilk mertebe bir iletim modelidir. Ortalama aktif gücü değiştirmeyen bir ters akım bile $I_{\mathrm{rms}}$ değerini ve kaybı büyütebilir. Nötr kol anahtarları, filtre endüktansları ve bağlantı iletkenleri aynı RMS artışından etkilenmektedir.

Kayıp dağılımının düşük güçteki değişimi de Eşitlik~\eqref{eq:iletim-kaybi-yaklasik} üzerinden yorumlanabilir. Düşük akımda $V_{\mathrm{on}}I_{\mathrm{avg,abs}}$ teriminin aktarılan çıkış gücüne oranı büyüdüğünden, yüksek güçte ikincil kalan sabit iletim gerilimi düşümü hafif yük veriminde belirginleşmektedir. Susturmalı ve Sin\textsuperscript{3} kompanzasyonlu çalışmada (S3) blok bazında ayrıştırılan kayıp, düşük akım bölgesinde $L_1$ çekirdek kaybının baskın olduğunu ve anahtarlama kaybının toplam kaybın yaklaşık \%7'sinde kaldığını göstermiştir (Alt Bölüm~\ref{sec:dusuk-akim-modelleme}). S1 ile S2 arasındaki kayıp farkı ise DA ve AA güçlerinin farkından değerlendirilmiştir.

Susturma, klasik PWM'nin bütün çalışma aralığında değiştirilmesi yerine belirlenen düşük güç bölgesinde etkinleştirilmelidir. Etkinleştirme kararı ölçülen akıma dayandırıldığında ters akımın güvenilir biçimde belirlenmesi, akım işaretinin yanında ek ölçütler gerektirmektedir. Ölçüm gürültüsü, akım sensörü ofseti, PWM güncelleme gecikmesi, ölü zaman ve anahtarlama anındaki diyot iletimi sıfır geçiş çevresinde yanlış karar üretebilir. Güvenilir uygulama akım eşiği, histerezis, minimum bekleme süresi ve güvenli komütasyon koşullarını birlikte kullanmalıdır. Akım denetleyicisinin sıfır geçiş çevresinde referans--ölçüm farkını azaltmak için oluşturduğu düzeltmeler, ana güç anahtarları ile nötr kol arasındaki geçiş sayısını da artırabilir. Susturmanın etkisi kalibre benzetimlerde kutup geriliminin ara seviyede kaldığı zaman oranı, DA ve AA güç farkı ve THD ile ölçülmüştür. $L_1$ akımı ve kapı işaretleri kayıtlarında ters akım aralığı da ölçülmüştür (Alt Bölüm~\ref{sec:s1-s2-kalibre}) \citep{Teodorescu2006PR,Zhang2014Deadtime}.

\section{Klasik Anahtarlama Yöntemi}

Klasik üç seviyeli modülasyonda nötr kol, sıfır gerilim durumunu üretmek üzere her iki yarım periyotta etkin kullanılmaktadır. Pozitif alternansta aktif enerji aktarımının yolu $P$ durumu, serbest dolaşımın yolu $O$ durumudur. Negatif alternansta aktif enerji aktarımı $N$ durumu, serbest dolaşım yine $O$ durumu üzerinden gerçekleşmektedir.

Benzetim ve deney, prototipin tek faz kolunu kapsamaktadır. Kapı işaretleri yarım periyot ayrımlı, tek kutuplu taşıyıcı karşılaştırmalı sinüzoidal PWM (SPWM) ile üretilmektedir (Alt Bölüm~\ref{sec:denetim-modelden}). Faz başına inceleme, nötr anahtar müdahalesinin hangi yarım periyotta ve hangi anahtar üzerinde yapıldığını açıkça göstermektedir.

Klasik yöntemde tamamlayıcı sürme, nötr kolu iki yönlü iletken yaparak akımın sürekliliğini korumaktadır. Düşük güçte aynı iki yönlü iletim dezavantaja dönüşmektedir. Akım sıfıra ulaştıktan sonra kapı komutunun korunması, akım yönünü sınırlayan doğal bir diyot kesimi bırakmamaktadır. Dolayısıyla klasik yöntem hafif yükte ters akım yükü ve RMS akım bakımından elverişsizdir. Düşük akımda ölü zamanın akım yönüne bağlı gerilim hatası da akım bozulmasını artırmaktadır (Alt Bölüm~\ref{sec:kalibrasyon}). İki yöntemin işlevsel farkı Tablo~\ref{tab:klasik-onerilen-islevsel} üzerinde özetlenmiştir.

\begin{table}[H]
  \centering
  \caption{Klasik ve önerilen anahtarlamanın düşük güçte işlevsel karşılaştırması.}
  \label{tab:klasik-onerilen-islevsel}
  \small
  \begin{tabularx}{\textwidth}{p{0.22\textwidth} X X}
    \toprule
    Özellik & Klasik senkron anahtarlama & Önerilen nötr anahtar susturma \\
    \midrule
    Sıfır durumu & Nötr kol iki yönlü akıma izin vermektedir. & Ters dolaşıma yol açan nötr komutu düşük güç bölgesinde engellenmektedir. \\
    Akım sürekliliği & Sıfır geçişten sonra ters yönde sürdürülebilir. & Akım sıfırda doğal olarak kesilebilir. Süreksiz iletim oluşabilir. \\
    Beklenen kayıp etkisi & Ters akım nedeniyle RMS ve iletim kaybı artabilir. & Ters akım yolu azaltıldığı için iletim kaybı düşebilir. \\
    Beklenen harmonik etkisi & Akım sürekliliği THD açısından avantaj sağlayabilir. & Üçüncü harmonik ve THD belirgin biçimde artabilir. \\
    Gerekli tamamlayıcı önlem & Özel düşük güç müdahalesi yoktur. & Sin\textsuperscript{3} kompanzasyonu, akım düzeyine bağlı uygulamada ise güvenli kip geçişi gerekmektedir. \\
    \bottomrule
  \end{tabularx}
\end{table}

\section{Önerilen Anahtarlama Yöntemi}

Önerilen yöntemin amacı, düşük güçte referans enerji yönüne ters akımı sürdüren nötr anahtar komutunu bastırmaktır. Yöntem dış anahtarların PWM yapısını değiştirmemekte, yalnız sıfır durumundaki karşı nötr anahtar komutunu kaldırmaktadır. Susturulan nötr anahtar, görev oranı komutu $u(t)$'nin işaretiyle, yani yarım periyotla belirlenmektedir.

Yarım periyot göstergesi $\kappa(t)\in\{0,1\}$
\begin{equation}
  \kappa(t)=
  \begin{cases}
    1, & u(t)\geq0,\\
    0, & u(t)<0
  \end{cases}
  \label{eq:susturma-etkinlestirme}
\end{equation}
ile tanımlansın. Klasik nötr kapı komutları $g_{Q3,\mathrm{kl}}(t)$ ve $g_{Q4,\mathrm{kl}}(t)$, susturmalı komutlar ise $g_{Q3,\mathrm{yeni}}(t)$ ve $g_{Q4,\mathrm{yeni}}(t)$ ile gösterildiğinde nötr anahtar susturma kuralı
\begin{equation}
  \begin{aligned}
    g_{Q3,\mathrm{yeni}}(t) &= g_{Q3,\mathrm{kl}}(t)\left[1-\kappa(t)\right],\\
    g_{Q4,\mathrm{yeni}}(t) &= g_{Q4,\mathrm{kl}}(t)\,\kappa(t)
  \end{aligned}
  \label{eq:notr-susturma-kurali}
\end{equation}
biçimindedir. Eşitlik~\eqref{eq:susturma-etkinlestirme} ve Eşitlik~\eqref{eq:notr-susturma-kurali} uyarınca pozitif yarım periyotta faz tarafı nötr anahtar $Q_3$, negatif yarım periyotta orta nokta tarafı nötr anahtar $Q_4$ susturulmaktadır. Dış anahtar komutları ve yük akımı yönündeki nötr anahtar komutu klasik yöntemdeki gibi kalmaktadır. İncelenen düşük güç çalışma noktasında susturma, benzetimde ve deneylerde senaryo boyunca sürekli etkindir. Susturmanın akım eşiği ve histerezisle yalnız düşük akımda etkinleştirilmesi, yöntemin akım düzeyine bağlı genişletmesidir. Uygulanan kapı mantığı aşağıda modelden okunmuştur.

\subsection{Modelden okunan kapı mantığı ve senaryo eşlemesi}
\label{sec:model-kapi-mantigi}

Senaryo tanımları ile benzetim modellerinin kapı üretim yapısı blok blok karşılaştırılmıştır. S1 modelinde dış anahtarlar pozitif ve negatif yarım periyotta PWM ile sürülmektedir. Nötr anahtar komutları ise PWM çıkışlarının mantıksal tümleyenidir. S2 ve S3 modellerinde dış anahtarların PWM üretimi aynıdır. Ancak nötr anahtarlar PWM'den bağımsız olarak yarım periyot boyunca sabit 1 veya 0 değerini almaktadır. S3'ün S2'den tek farkı, modülasyon referansına eklenen Sin\textsuperscript{3} terimidir. Tablo~\ref{tab:model-kapi-mantigi} iki kapı mantığını ideal durumlar üzerinden karşılaştırmaktadır. Dört kapı komutunun her biri \SI{1}{\micro\second} ölü zamanla yükselen kenarda geciktirilmektedir. Hiçbir kapı çiftinde örtüşme yoktur.

\begin{table}[H]
  \centering
  \caption{Benzetim modellerinden okunan kapı mantığı: klasik anahtarlama (S1) ve nötr anahtar susturma (S2, S3).}
  \label{tab:model-kapi-mantigi}
  \small
  \begin{tabular}{l c c c c c c c c c}
    \toprule
    & & \multicolumn{4}{c}{S1 (klasik)} & \multicolumn{4}{c}{S2, S3 (susturma)} \\
    \cmidrule(lr){3-6}\cmidrule(lr){7-10}
    Yarım periyot & Durum & $Q_1$ & $Q_2$ & $Q_3$ & $Q_4$ & $Q_1$ & $Q_2$ & $Q_3$ & $Q_4$ \\
    \midrule
    $u\geq0$ & $P$ & 1 & 0 & 0 & 1 & 1 & 0 & 0 & 1 \\
    $u\geq0$ & $O$ & 0 & 0 & \textbf{1} & 1 & 0 & 0 & \textbf{0} & 1 \\
    $u<0$ & $N$ & 0 & 1 & 1 & 0 & 0 & 1 & 1 & 0 \\
    $u<0$ & $O$ & 0 & 0 & 1 & \textbf{1} & 0 & 0 & 1 & \textbf{0} \\
    \bottomrule
  \end{tabular}
\end{table}

Tablo~\ref{tab:model-kapi-mantigi} içindeki kalın hücreler iki yöntem arasındaki tek farkı göstermektedir. Klasik anahtarlamada $O$ durumunda karşı nötr anahtar da iletimdedir ve ters akım yolu açıktır. Susturmalı anahtarlamada $O$ durumunda yalnız yük akımı yönündeki yol açıktır. Pozitif yarım periyotta akım orta noktadan $Q_4$ ve $DS_3$ üzerinden faza akmaktadır. Ters yönde ise akım yolu yoktur. $L_1$ akımı sıfıra indiğinde kesilmekte ve kutup gerilimi filtre tarafının gerilimini izlemektedir. $D_1$--$D_2$ yalnız kutup gerilimi DA rayını aşarsa iletime girmektedir. S1 kapı mantığı Tablo~\ref{tab:t-tipi-durumlar} ile birebir aynıdır.

Susturmanın devre düzeyindeki sonucu Şekil~\ref{fig:susturma-esdeger} üzerinde gösterilmiştir. Klasik anahtarlamada her yarım periyot, iki yönlü iletebilen alt ve üst anahtardan oluşan senkron alçaltıcıya karşılık gelmektedir (Şekil~\ref{fig:senkron-buck-esdeger}). Susturmalı çalışmada iletimde kalan nötr anahtar ile karşı nötr anahtarın antiparalel diyodu seri bir tek yönlü eleman oluşturmaktadır. Faz kolunun her yarım periyottaki eşdeğeri, bir aktif anahtar ile bir diyottan oluşan senkron olmayan alçaltıcıdır. Serbest dolaşım akımı diyot üzerinden yalnız yük yönünde akmakta ve sıfıra indiğinde kendiliğinden kesilmektedir. Ters akımın kesilmesi ile süreksiz iletimin başlaması aynı devre özelliğinin iki sonucudur.

\begin{figure}[H]
  \centering
  \begin{subfigure}[t]{0.48\textwidth}
    \centering
    \resizebox{\linewidth}{!}{\susturmabuckvector{$+V_{\mathrm{dc}}/2$}{$N$}{$Q_1$}{$Q_4$ iletimde\\$+\,DS_3$}{1}}
    \caption{Pozitif yarım periyot: $Q_1$ anahtarlanmakta, $Q_3$ susturulmaktadır}
  \end{subfigure}\hfill
  \begin{subfigure}[t]{0.48\textwidth}
    \centering
    \resizebox{\linewidth}{!}{\susturmabuckvector{$N$}{$-V_{\mathrm{dc}}/2$}{$Q_2$}{$Q_3$ iletimde\\$+\,DS_4$}{0}}
    \caption{Negatif yarım periyot: $Q_2$ anahtarlanmakta, $Q_4$ susturulmaktadır}
  \end{subfigure}
  \caption{Nötr anahtar susturmada faz kolunun yarım periyot eşdeğerleri: senkron olmayan iki alçaltıcı dönüştürücü. Mavi yol aktif enerji aktarımını, turuncu kesikli yol diyot üzerinden tek yönlü serbest dolaşımı göstermektedir. Ters yöndeki akım için nötr koldan yol yoktur.}
  \label{fig:susturma-esdeger}
\end{figure}

Benzetim modelinde S2 kapı yolu, tezdeki nötr anahtar susturma tanımına göre kurulmuştur (Alt Bölüm~\ref{sec:model-dogrulama-duzeltme}). Kapı yolunun dört komutu Şekil~\ref{fig:s2-kapi-isaretleri} üzerinde bir şebeke periyodu ve dört anahtarlama periyodu boyunca gösterilmiştir. $Q_1$ pozitif yarım periyotta \SI{10}{\kilo\hertz} PWM ile anahtarlanırken $Q_2$ kesimde, $Q_4$ iletimde ve $Q_3$ kesimdedir. Negatif yarım periyotta roller simetrik olarak yer değiştirmektedir. Nötr anahtarlar yarım periyot boyunca anahtarlama yapmadan sabit konumda kalmaktadır. Dolayısıyla anahtarlama kaybı yalnız dış anahtarlarda oluşmaktadır.

\begin{figure}[H]
  \centering
  \includegraphics[width=0.96\textwidth]{model_dogrulama/s2_kapi_isaretleri.pdf}
  \caption[Nötr anahtar susturmada (S2) $Q_1$--$Q_4$ kapı komutları]{Nötr anahtar susturmada (S2) $Q_1$--$Q_4$ kapı komutları: solda son tam şebeke periyodu (\num{80}--\SI{100}{\milli\second}), sağda pozitif yarım periyottan \SI{400}{\micro\second} kesit. Kayıt, doğrulanmış benzetim modelinden alınmıştır. Nötr anahtarlar yarım periyot boyunca sabit 1 veya 0 değerindedir.}
  \label{fig:s2-kapi-isaretleri}
\end{figure}

Klasik anahtarlamanın aynı pencerelerdeki kapı komutları Şekil~\ref{fig:s1-kapi-isaretleri} üzerinde verilmiştir. Pozitif yarım periyotta $Q_1$ ile $Q_3$ tamamlayıcı darbelerle anahtarlanırken $Q_2$ kesimde, $Q_4$ iletimde tutulmaktadır. Negatif yarım periyotta $Q_2$ ile $Q_4$ tamamlayıcı çalışmaktadır. İki şekil karşılaştırıldığında tek fark, S2'de sıfır durumundaki karşı nötr anahtar komutunun ($Q_3$ pozitif, $Q_4$ negatif yarım periyotta) bastırılmasıdır. Kapı dizileri Tablo~\ref{tab:model-kapi-mantigi} ile örtüşmektedir. Ölü zaman ve kapı polaritesi model bloklarından okunmuştur.

\begin{figure}[H]
  \centering
  \includegraphics[width=0.96\textwidth]{model_dogrulama/s1_kapi_isaretleri.pdf}
  \caption[Klasik anahtarlamada (S1) $Q_1$--$Q_4$ kapı komutları]{Klasik tamamlayıcı anahtarlamada (S1) $Q_1$--$Q_4$ kapı komutları: solda son tam şebeke periyodu (\num{80}--\SI{100}{\milli\second}), sağda pozitif yarım periyottan \SI{400}{\micro\second} kesit. Kayıt, kalibre S1 benzetim modelinden alınmıştır.}
  \label{fig:s1-kapi-isaretleri}
\end{figure}

Susturmanın donanımda güvenli uygulanması için dört koşul birlikte sağlanmalıdır. İlk olarak, akım düzeyine bağlı etkinleştirmede akım işareti sensör ofsetinden büyük bir güven payıyla belirlenmelidir. İkinci olarak, nötr kolun kesilmesi endüktör akımı için kesintisiz ve aşırı gerilim üretmeyen alternatif bir yol bırakmalıdır. Susturmalı çalışmada akım sıfıra inince yol kendiliğinden kapanmaktadır. Ters akım yönünde ise $D_1$--$D_2$ ve dış anahtar diyotları yol sağlamaktadır. Üçüncü olarak, karşılıklı anahtarların ölü zamanı ve minimum darbe genişliği korunmalıdır. Benzetim modelinde ölü zaman $t_d=\SI{1}{\micro\second}$'dir. Değer, donanımın PWM zamanlayıcısındaki ölü zaman yazmacından gelen tasarım değeridir. Düşük akımda ölü zamanın akıma bağlı gerilim hatası Alt Bölüm~\ref{sec:kalibrasyon} içinde değerlendirilmiştir \citep{Zhang2014Deadtime}. Son olarak, ölçüm geçersizliği veya aşırı akım durumunda denetim önceden tanımlanmış güvenli anahtarlamaya dönmelidir.

Nötr anahtarın bastırılması ters akım yolunu azaltırken endüktör akımını süreksiz iletime zorlamaktadır. Akımın sıfır çevresindeki düzleşmesi ve iletim aralıklarının daralması, temel bileşene göre üçüncü harmonik payını büyüten etkenlerdir. Süreksiz iletimde görev oranı ile anahtarlama anındaki endüktör gerilim farkı birlikte değiştiğinden akım zarfı saf sinüs biçiminden ayrılmaktadır. Alternans ortasında ($\omega_0t=\pi/2$ ve $3\pi/2$ çevresinde) akım zarfındaki çökmenin nedeni, DA bara ile şebeke gerilimi arasındaki akım yükseltme payının azalmasıdır. Sıfır geçişlerine yakın bölgelerdeki farklı yükselme hızı da üçüncü harmonik bileşeni beslemektedir. Harmonik cezanın azaltılması için Bölüm~\ref{chap:sin3} altında incelenen Sin\textsuperscript{3} ileri besleme kompanzasyonu, önerilen anahtarlamayı tamamlayan katman olarak kullanılmaktadır.

\section{Verim ve Harmonik Performans Ödünleşimi}

Deney senaryoları aynı düşük güç probleminin ardışık karşılaştırma kollarıdır: S1 klasik senkron anahtarlama, S2 nötr anahtar susturma, S3 ise susturma ile Sin\textsuperscript{3} kompanzasyonunun birlikte kullanımıdır. Senaryo eşlemesi benzetim modellerinin kapı mantığından doğrulanmıştır (Alt Bölüm~\ref{sec:model-kapi-mantigi}). S1--S3 senaryolarında modülasyon referansı, akım denetleyicisi olmadan açık çevrimde uygulanmıştır. Her senaryo kendi deneyinin DA bara ve şebeke gerilimiyle benzetilmiştir. Modelde senaryolar arasında nötr kapı mantığı, modülasyon genliği ve referans fazı değişmiştir. Açık çevrim koşulu yöntemin oluşum mekanizmasını ayırmak için yararlıdır. Kapalı çevrim performansı Alt Bölüm~\ref{sec:s4-kapali-cevrim} içinde ayrıca ele alınmıştır.

\subsection{Kalibre benzetimde klasik anahtarlama ile susturmanın karşılaştırılması}
\label{sec:s1-s2-kalibre}

S1 ile S2, deneyin aktif ve reaktif gücüne kalibre edilmiş benzetim modeliyle karşılaştırılmaktadır. Kalibrasyon yöntemi, kalibre noktalar ve S3 dahil bütün deney ve benzetim değerleri Alt Bölüm~\ref{sec:kalibrasyon} içinde verilmiştir. Aktif güç ve reaktif güç kalibrasyon hedefidir. Şebeke akımı, verim, DA güç ve THD ise bağımsız doğrulama büyüklükleridir.

Tablo~\ref{tab:kalibrasyon-karsilastirma} içindeki S1 ve S2 sütunları üç bulguyu desteklemektedir. Birinci bulgu verimdir. Susturma, kalibre benzetimde verimi \%91,15'ten \%93,93'e, yani 2,78 puan artırmaktadır. Deneyde artış \%92,5'ten \%94,2'ye, 1,7 puandır. Verim sıralaması benzetim ile deneyde aynıdır. S2'de model DA gücünü \%0,004, verimi $-0{,}27$ puan hatayla yeniden üretmektedir. S1'de benzetim kaybı \SI{24.0}{\watt}, deney kaybı \SI{19.7}{\watt}'tır ve S1 kaybı diyot toparlanma parametrelerine duyarlıdır. Verim farkının belirsizliği, 0,1 saniyedeki oturma farkı nedeniyle S2 için yaklaşık 0,1 puandır.

İkinci bulgu harmonik cezadır. Susturma, şebeke (faz) akımı THD'sini deneyde \%15'ten \%47'ye, kalibre benzetimde \%39,5'ten \%56,7'ye yükseltmektedir. Artış iki kaynakta aynı yöndedir ve iki kaynakta da S2 en yüksek THD'li senaryodur. Deney değerleri AA analizörünün kayıt süresi boyunca izlenen THDi değerleridir ve osiloskop kayıtlarından hesaplanan değerlerle uyumludur (Tablo~\ref{tab:thd-kaynaklar}). Model THD sıralamasını deneyle aynı vermekte, S3'te deneye çok yakın değer üretmektedir. S1 ve S2'deki mutlak değer farkları, dalga biçimi karşılaştırması ve olası kaynaklarıyla birlikte kalibrasyon incelemesinde ele alınmıştır.

Üçüncü bulgu, susturmanın köprü davranışını değiştirmesidir. Kutup geriliminin ara seviyede kaldığı zaman oranı S1'de \%0,5, S2'de \%12,4'tür. Susturmalı köprü düşük akımda akım kaynağına yaklaşmakta ve aktif güç modülasyon genliğine daha duyarlı hâle gelmektedir (Tablo~\ref{tab:kopru-ic-empedansi}). Kalibre S2 genliğinin parametre dosyasındaki ilk değerden \%3,3 düşük olması da aynı davranışın sonucudur.

Şekil~\ref{fig:s1-s2-kalibre-dalga}, iki kalibre noktanın son şebeke periyodundaki dalga biçimlerini karşılaştırmaktadır. S1'de şebeke akımı sinüs zarfının üzerinde anahtarlama kaynaklı bozulma taşımaktadır. S2'de ise akım zarfı alternans ortasında çökmekte ve sıfır geçişlerinin iki yanında tepe yapmaktadır. Akım zarfının biçimi, Şekil~\ref{fig:susturma-esdeger} ile açıklanan süreksiz iletimle ve üçüncü harmonik artışıyla uyumludur. Alt panel aynı zaman aralığında iki senaryonun kutup gerilimini göstermektedir. S2'de sıfır durumundan sonra $L_1$ akımı söndüğünde kutup gerilimi ara bir seviyeye sıçramakta ve bir sonraki darbeye kadar ara seviyede kalmaktadır. Şekil~\ref{fig:s1-s2-kalibre-l1}, aynı kalibre noktalarda $L_1$ akımını göstermektedir. Kayıt, kalibrasyon değerleri değiştirilmeden $L_1$ akımı ve kapı işaretleri de kaydedilerek yinelenen benzetimlerden alınmıştır. S1'de yaklaşık \SI{9}{\ampere} tepe--tepe dalgalanma temel bileşenden büyük olduğundan $L_1$ akımı her anahtarlama periyodunda işaret değiştirmektedir. Pozitif yarım periyodun \%25,2'sinde akım $-\SI{0.05}{\ampere}$'in altındadır, yani ters yöndedir. 400~\si{\micro\second} kesitte akım her periyotta yaklaşık $-\SI{2.6}{\ampere}$'e inmektedir. S2'de ters akım aralığı yoktur. Akım sıfıra indiğinde bir sonraki darbeye kadar sıfırda kalmaktadır. Pozitif yarım periyodun \%17'sinde akım sıfırdır. Akım kaydı, süreksiz iletim açıklamasının ve susturmanın ters akım yolunu kestiğinin doğrudan görsel kanıtıdır.

\begin{figure}[H]
  \centering
  \includegraphics[width=0.96\textwidth]{model_dogrulama/s1_s2_kalibre_l1.pdf}
  \caption[Kalibre benzetimde S1 ve S2 $L_1$ akımı]{Kalibre benzetimde $L_1$ akımı: (a) S1 ve (c) S2 için son tam şebeke periyodunda $i_{L1}$ ve şebeke akımı $i_g$, (b) S1 ve (d) S2 için pozitif yarım periyottan \SI{400}{\micro\second} kesit. S1'de akım her anahtarlama periyodunda ters yöne geçmekte, S2'de sıfırda kalmaktadır.}
  \label{fig:s1-s2-kalibre-l1}
\end{figure}

\begin{figure}[H]
  \centering
  \includegraphics[width=0.96\textwidth]{model_dogrulama/s1_s2_kalibre_dalga.pdf}
  \caption[Kalibre benzetimde S1 ve S2 kutup gerilimi, şebeke gerilimi ve şebeke akımı]{Kalibre benzetimde son tam şebeke periyodu (\num{80}--\SI{100}{\milli\second}): (a) S1 ve (b) S2 için kutup gerilimi $v_{\mathrm{inv}}$, C noktasındaki gerilim $v_g$ ve şebeke akımı $i_g$. (c) Pozitif yarım periyottan \SI{400}{\micro\second} kesitte iki senaryonun kutup gerilimi. S2'de akım söndüğünde kutup gerilimi ara bir seviyeye geçmektedir.}
  \label{fig:s1-s2-kalibre-dalga}
\end{figure}

\section{Bölüm Sonucu}

T-tipi faz kolunun düşük güçte senkron buck benzeri davranışı, nötr kol üzerinden ters enerji dolaşımının fiziksel kaynağını açıklamaktadır. Deney çalışma noktasının tepe akımı yaklaşık \SI{9}{\ampere} tepe--tepe $L_1$ dalgalanmasının yarısından küçük olduğundan klasik anahtarlama temel periyodun büyük bölümünde ters akıma izin vermektedir. Önerilen anahtarlama, ters akımı sürdüren nötr komutunu düşük güç bölgesinde bastırarak kayıp yolunu azaltmaktadır. Yöntemin bedeli ise süreksiz iletim ve üçüncü harmonik cezasıdır. Modelden okunan kapı mantığı, iki yöntem arasındaki tek farkın $O$ durumundaki karşı nötr anahtar komutu olduğunu doğrulamıştır.

Susturma deneyde verimi \%92,5'ten \%94,2'ye yükseltmiş, akım THD'sini ise \%15'ten \%47'ye çıkarmıştır. Harmonik cezayı azaltma gereksinimi, önerilen anahtarlamayı tamamlayan Sin\textsuperscript{3} kompanzasyonunun gerekçesidir.
